\documentclass[a4paper, 12pt]{report}

\usepackage[spanish]{babel}
\usepackage[utf8]{inputenc}
\usepackage[left=4cm, right=4cm, top=4cm]{geometry}
\usepackage{textcomp}
\usepackage{booktabs}
\usepackage{amssymb}
\usepackage{bussproofs}
\usepackage{fancyhdr}
\usepackage{graphicx}
\usepackage{amsmath}
\usepackage{enumitem}
\usepackage{ifsym}


\usepackage{hyperref}
\hypersetup{
    colorlinks=true,
    linkcolor=blue,
    filecolor=magenta,
    urlcolor=cyan,
}

\pagestyle{fancy}
\lhead{Almeida, Figueroa \& Ibarra}
\chead{Tarea 4}
\rhead{\today}

\begin{document}
\begin{titlepage}
    \centering
    {\scshape\Huge Universidad Nacional Autónoma de México \par}
    \vspace{1.25cm}
    {\scshape\huge Fundamentos de Bases de Datos\par}
    \vspace{1.25cm}
    {\huge\bfseries Tarea 4: Álgebra Relacional\par}
    \vspace{1.25cm}
    {\Large\textsc Almeida Rodríguez Jerónimo\par}
    \vspace{.1cm}
    {\large\texttt{418003815}\par}
    \vspace{0.25cm}
    {\Large\textsc Figueroa Sandoval Gerardo Emiliano\par}
    \vspace{.1cm}
    {\large\texttt{315241774}\par}
    \vspace{0.25cm}
    {\Large\textsc Ibarra Moreno Gisselle \par}
    \vspace{.1cm}
    {\large\texttt{315602193}\par}
    \vspace{1.5cm}
    \vfill
    \begin{figure}[hb!]
        \includegraphics[width=.3\textwidth]
            {../logos/escudo_f-ciencias.png}\hfill
        \includegraphics[width=.3\textwidth]
            {../logos/Escudo_UNAM.png}\hfill
    \end{figure}
\end{titlepage}

\section*{Ejercicio 1}{
\begin{enumerate}[label=\alph*)]
\item{Jero}
\item{ Y user; count(id\_fotografía) $\rightarrow$ total\_fotos (Subir)\\
}
\item{}
\item{Jero}
\item{r = Y id\_fotografía; count(user) $\rightarrow$ megustas (calificar)\\
s = $\pi$ user, id\_fotografía, megustas (Subir $\Join$ r )\\
t = Y user; max(megustas) $\rightarrow$ max\_megustas (s)\\
$\pi$ user, id\_fotografía (s $\Join$ t )}
\end{enumerate}
}

\section*{Ejercicio 2}{
\begin{enumerate}[label=\alph*)]
\item{}
\item{¿Qué fabricantes producen computadoras portátiles con un disco duro de
    menos 100 GB?}
\item{r = $\sigma$ fabricante = 'B' (Producto)\\
	s = $\pi$ modelo, precio (Laptop) $\cup$ $\pi$ modelo, precio (PC)
	$\cup$ $\pi$ modelo, precio (Impresora) \\
	$\pi$ modelo, precio (s $\Join$ r)}
\item{}
\item{Encontrar los números de modelo de todas las impresoras láser a color.}
\item{r = $\pi$ modelo, fabricante (Producto)\\
	s = $\pi$ fabricante ($\pi$ modelo (Laptop) $\Join$ r)\\
	t = $\pi$ fabricante ($\pi$ modelo (PC) $\Join$ r)\\
	s - t\\}
\item{}
\item{Encontrar toda la información de las PCs que tienen la misma velocidad y
    RAM.}
\item{r = $\pi$ modelo ($\sigma$ velocidad $\geq$ 2.8 (PC))\\
s = $\pi$ modelo($\sigma$ velocidad $\geq$ 2.8 (Laptop))\\
$\pi$ fabricante ((r $\cup$ s) $\Join$ Producto)\\
}
\item{}
\item{Jero}
\item{r = $\pi$ modelo, fabricante (Producto $\Join$ PC)\\
s = Y fabricante; count(modelo) $\rightarrow$ numproductos (r)\\
$\pi$ fabricante ($\sigma$ numproductos = 3 (s))}
\item{}
\item{Jero}
\item{r = $\pi$ modelo ($\sigma$ fabricante = 'E' (Producto)) $\Join$ Laptop\\
s = $\sigma$ hd $<$ 200 (r)\\
$\pi$ modelo, velocidad, ram, hd\_nuevo $\leftarrow$ hd * 1.15, pantalla, precio (s)\\
}
\item{}
\item{Jero}
\item{}
\item{$f(x)=xe^{-2x}\sin(x)$\\
           Proponemos la siguiente solución:\\
           \begin{align*}
               y_p&=(Ax+B)e^{-2x}\sin(x)+(Cx+D)e^{-2x}\cos(x)\\
               y_p'&=xe^{-2x}\sin(x)-2(Ax+B)e^{-2x}\sin(x)+(Ax+B)e^{-2x}\cos(x)\\
                   &+xe^{-2x}\cos(x)-2(Cx+D)e^{-2x}\cos(x)-(Cx+D)e^{-2x}\sin(x)\\
                   &=xe^{-2x}\big[\sin(x)+\cos(x)\big]
                        +(Ax+B)e^{-2x}\big[-2\sin(x)+\cos(x)\big]\\
                   &-(Cx+D)e^{-2x}\big[2\cos(x)+\sin(x)\big]\\
               y_p''&=-2e^{-2x}\Big[(-2Ax+A-2B)[C\sin(x)+D\cos(x)]+(Ax+B)[C\cos(x)-D\sin(x)]\Big]\\
                   &+e^{-2x}\Big[-2Ax[C\sin(x)+D\cos(x)]+(-2Ax+A-2B)[C\cos(x)-D\sin(x)]\\
                   &+Ax[C\cos(x)-D\sin(x)]+(Ax+B)[-C\sin(x)-D\cos(x)]\Big]\\
                   &=e^{-2x}\Big[(Ax-2A+3B)[C\sin(x)+D\cos(x)]+(A-B)[C\cos(x)-D\sin(x)]\Big]
           \end{align*}
           Entonces,
               Entonces,
       \begin{align*}
           xe^{-2x}\sin(x)&=e^{-2x}\big[(Ax-2A+3B)[C\sin(x)+D\cos(x)]+(A-B)[C\cos(x)-D\sin(x)]\big]\\
           &+2e^{-2x}\big[(-2Ax+A-2B)[C\sin(x)+D\cos(x)]+(Ax+B)[C\cos(x)-D\sin(x)]\big]\\
           &-3(Ax+B)e^{-2x}[C\sin(x)+D\cos(x)]\\
           x\sin(x)&=\big[(Ax-2A+3B)[C\sin(x)+D\cos(x)]+(A-B)[C\cos(x)-D\sin(x)]\big]\\
           &+\big[(-4Ax+2A-4B)[C\sin(x)+D\cos(x)]+(2Ax+2B)[C\cos(x)-D\sin(x)]\big]\\
           &-(3Ax+3B)[C\sin(x)+D\cos(x)]\\
           x\sin(x)&=\big[(-6Ax-4B)[C\sin(x)+D\cos(x)]+(2Ax+A+B)[C\cos(x)-D\sin(x)]\big]\\
                   &=(-2Ax(3C+D)-B(4C+D)+AD)\sin(x)\\
                   &+(-2Ax(3D+C)-B(4D+C)+AC)\cos(x)\\
       \end{align*}
       }
\end{enumerate}
}



\end{document}
